% ===== PRÉAMBULE CANONIQUE — à coller VERBATIM en tête de CHAQUE .tex =====
\documentclass[11pt,a4paper]{article}
\usepackage[utf8]{inputenc}\usepackage[T1]{fontenc}\usepackage{lmodern}
\usepackage[french]{babel}
\usepackage{amsmath,amssymb,mathtools}\usepackage{icomma}
\usepackage[version=4]{mhchem}          % \ce{...} : équations & formules
\usepackage{siunitx}\sisetup{locale=FR} % \qty{}{}, \si{}, \ang{}
\usepackage{chemfig}                    % \chemfig{...} : structures développées
\usepackage{xcolor}\usepackage{colortbl}
\usepackage{tikz}\usepackage{pgfplots}\pgfplotsset{compat=1.18}
\usepackage[most]{tcolorbox}\usepackage{enumitem}\usepackage{array,booktabs}
\usepackage[a4paper,margin=2.2cm]{geometry}
\usetikzlibrary{arrows.meta,calc,angles,quotes,positioning,patterns,backgrounds,shapes.geometric}
\definecolor{primary}{RGB}{222,49,99}\definecolor{primarydark}{RGB}{150,22,55}
\definecolor{defcol}{RGB}{0,114,178}\definecolor{methcol}{RGB}{52,92,148}
\definecolor{excol}{RGB}{0,135,90}\definecolor{attncol}{RGB}{200,40,55}
\tcbset{boxbase/.style={enhanced,breakable,boxrule=0.9pt,arc=2.5pt,
  left=3mm,right=3mm,top=2.3mm,bottom=2.3mm,fonttitle=\bfseries,coltitle=white}}
% --- boîtes TUTORIEL (noms EXACTS, ne pas renommer) ---
\newtcolorbox{cle}[1][]{boxbase,colback=defcol!6,colframe=defcol,
  title={\textsc{À retenir}\ifx\relax#1\relax\relax\else\textmd{ : \itshape #1}\fi}}
\newtcolorbox{methode}[1][]{boxbase,colback=methcol!7,colframe=methcol,sharp corners,
  title={\textsc{Méthode}\ifx\relax#1\relax\relax\else\textmd{ --- \itshape #1}\fi}}
\newtcolorbox{exemplebox}[1][]{boxbase,colback=excol!5,colframe=excol,
  title={\textsc{Exemple}\ifx\relax#1\relax\relax\else\textmd{ : \itshape #1}\fi}}
\newtcolorbox{piege}[1][]{boxbase,colback=attncol!6,colframe=attncol,
  title={\textsc{Piège}\ifx\relax#1\relax\relax\else\textmd{ : \itshape #1}\fi}}
\newtcolorbox{remarque}[1][]{boxbase,colback=black!4,colframe=black!45,
  title={\textsc{Remarque}\ifx\relax#1\relax\relax\else\textmd{ : \itshape #1}\fi}}
% --- boîtes EXERCICE / CORRIGÉ (noms EXACTS, auto-comptées) ---
\newtcolorbox[auto counter]{exo}[1][]{boxbase,colback=primary!4,colframe=primary,
  title={\textsc{Exercice}~\thetcbcounter\ifx\relax#1\relax\relax\else\textmd{ --- \itshape #1}\fi}}
\newtcolorbox[auto counter]{corrige}[1][]{boxbase,colback=excol!4,colframe=excol,
  title={\textsc{Corrigé}~\thetcbcounter\ifx\relax#1\relax\relax\else\textmd{ --- \itshape #1}\fi}}
\newcommand{\R}{\mathbb{R}}\newcommand{\N}{\mathbb{N}}\newcommand{\Z}{\mathbb{Z}}\newcommand{\ds}{\displaystyle}
% =========================================================================
\begin{document}
% THEME_TITLE: Équilibre chimique et loi de Le Chatelier
\begin{center}
{\Large\bfseries\color{primarydark} Thème 5 — Équilibre chimique et Le Chatelier}\\[2pt]
{\color{primary}\small Constante $K$, loi d'action de masse, déplacement d'équilibre}
\end{center}

\begin{cle}[l'état d'équilibre et la constante $K$]
Une réaction \textbf{réversible} ($\rightleftharpoons$) atteint un \textbf{équilibre} quand la vitesse du sens direct égale celle du sens inverse : les concentrations \emph{ne varient plus} (état stationnaire), mais les deux réactions continuent (\textbf{équilibre dynamique}). Pour $a\,\ce{A} + b\,\ce{B} \rightleftharpoons c\,\ce{C} + d\,\ce{D}$ :
\[ \boxed{\;K = \dfrac{[\ce{C}]^{c}\,[\ce{D}]^{d}}{[\ce{A}]^{a}\,[\ce{B}]^{b}}\;} \qquad\text{(produits sur réactifs, coefficients en exposants).} \]
$K$ est \textbf{sans unité} et ne dépend que de la \textbf{température}. Interprétation : $K\gg 1 \Rightarrow$ équilibre vers les \textbf{produits} ; $K\ll 1 \Rightarrow$ vers les \textbf{réactifs} ; $K\approx 1 \Rightarrow$ coexistence.
\end{cle}

\begin{methode}[écrire $K$ et exploiter la loi de Le Chatelier]
\textbf{Écrire $K$} : produits au numérateur, réactifs au dénominateur, chaque exposant $=$ coefficient. \textbf{Les solides purs et liquides purs n'apparaissent pas} (seuls les gaz et les espèces $(aq)$ figurent).\\[3pt]
\textbf{Le Chatelier} : perturbé, le système évolue dans le sens qui \emph{s'oppose} à la perturbation.
\renewcommand{\arraystretch}{1.25}
\begin{center}\small
\begin{tabular}{@{}>{\raggedright\arraybackslash}p{4.0cm}>{\raggedright\arraybackslash}p{9.3cm}@{}}
\toprule
\textbf{Perturbation} & \textbf{Sens du déplacement}\\
\midrule
$T \nearrow$ & vers le sens \textbf{endothermique} (qui absorbe la chaleur)\\
$p \nearrow$ (gaz) & vers le côté qui compte \textbf{moins} de moles de gaz\\
ajout d'un réactif & sens qui \textbf{consomme} ce réactif (vers les produits)\\
catalyseur ; solide/liquide pur & \textbf{aucun déplacement}\\
\bottomrule
\end{tabular}
\end{center}
\renewcommand{\arraystretch}{1.0}
Seule la \textbf{température} modifie la \emph{valeur} de $K$.
\end{methode}

\begin{exemplebox}[calculer $K$ à partir des concentrations]
Pour $\ce{PCl5}(g) \rightleftharpoons \ce{PCl3}(g) + \ce{Cl2}(g)$ à l'équilibre, avec $[\ce{PCl5}]=\qty{0.40}{\mol\per\litre}$, $[\ce{PCl3}]=\qty{1.40}{\mol\per\litre}$, $[\ce{Cl2}]=\qty{0.50}{\mol\per\litre}$ :
\[ K = \frac{[\ce{PCl3}]\,[\ce{Cl2}]}{[\ce{PCl5}]} = \frac{1{,}40\times 0{,}50}{0{,}40} = 1{,}75\;(>1). \]
\end{exemplebox}

\begin{piege}[les fautes classiques]
\begin{itemize}[leftmargin=1.6em,topsep=2pt,itemsep=1.5pt]
  \item ne \emph{jamais} mettre un solide pur ou un liquide pur dans $K$.
  \item «\,équilibre\,» $\ne$ «\,quantités égales\,» : les concentrations sont \emph{constantes}, pas forcément égales.
  \item un \textbf{catalyseur} ne déplace pas l'équilibre (il n'agit que sur la vitesse).
  \item l'effet de la \textbf{pression} n'existe que si le nombre de moles de gaz \emph{change} entre les deux côtés ; multiplier les coefficients par $n$ élève $K$ à la puissance $n$.
\end{itemize}
\end{piege}

\end{document}
